% Zumdahl Chapter 10 free-response set -- 2 long (10) + 3 short (4) = 32 pts.
\documentclass{apchem-exam}

\apcunitword{Chapter}
\apcunit{10}
\apcunittitle{Liquids and Solids}
\apcdoctitle{Free-Response Set}
\apcsubtitle{Zumdahl \S10.1--\S10.8 \textbullet\ 5 questions \textbullet\ 32 points \textbullet\ 50 minutes}

\begin{document}
\apctitleblock

\begin{apcnote}[Scope --- one CED exclusion applies]
Covers \textbf{CED 3.1} (intermolecular forces, \S10.1), \textbf{3.2}
(properties of solids, \S10.3--10.7), \textbf{3.3} (\S10.2, \S10.8),
\textbf{2.3} (ionic solids, \S10.7), \textbf{2.4} (metals and alloys,
\S10.4) and \textbf{6.5} (energy of phase changes, \S10.8).

\textbf{Phase diagrams (\S10.9) are excluded} by the CED statement on
topic 3.3 and are not examined here. Unit-cell geometry (\S10.7) is
enrichment depth under the exclusion on 2.3: the \emph{types} of solid and
their properties are assessed, specific crystal structures are not.

Take $c(\ce{H2O},\ell) = \SI{4.18}{\joule\per\gram\per\celsius}$,
$\Delta H_{\text{fus}} = \SI{334}{\joule\per\gram}$ and
$\Delta H_{\text{vap}} = \SI{2260}{\joule\per\gram}$.
\end{apcnote}

\examsection{II}{Free Response}{5 questions \textbullet\ 32 points \textbullet\ 50 minutes}{\calculatorok}

% =====================================================================
\frq{long}{10}{Zumdahl \S10.1 \textbullet\ CED 3.1 \textbullet\ intermolecular forces and boiling points}

The normal boiling points of four substances are:
\begin{center}
\begin{tabular}{lc}
\hline
\textbf{Substance} & \textbf{Boiling point (\si{\celsius})} \\
\hline
\ce{CH4}  & $-164$ \\
\ce{HCl}  & $-85$ \\
\ce{HF}   & $+20$ \\
\ce{H2O}  & $+100$ \\
\hline
\end{tabular}
\end{center}

\begin{parts}
  \item Identify the strongest intermolecular force present in each
        substance. \answerlines[3]{\ce{CH4}: London dispersion only. \quad
        \ce{HCl}: dipole--dipole. \quad \ce{HF}: hydrogen bonding. \quad
        \ce{H2O}: hydrogen bonding.}
  \item Explain why \ce{HF} boils higher than \ce{HCl} even though
        \ce{HCl} has more electrons. \workspace[3.2cm]
        \keyonly{\begin{apcanswer}More electrons would give \ce{HCl} the
        stronger dispersion forces, so on that basis alone it should boil
        higher. It does not, because a different and much stronger force
        operates in \ce{HF}.

        Fluorine is small and highly electronegative, so the \ce{H-F} bond
        is strongly polarized and the hydrogen is left with very little
        electron density. It can approach a lone pair on a neighbouring
        fluorine very closely, and the resulting hydrogen bond is far
        stronger than the dipole--dipole attraction in \ce{HCl}. Chlorine
        is larger and less electronegative, so \ce{HCl} cannot hydrogen
        bond.\end{apcanswer}}
  \item Explain why \ce{H2O} boils higher than \ce{HF}, even though
        fluorine is more electronegative than oxygen.
        \answerlines[3]{Each water molecule has \emph{two} hydrogens and
        \emph{two} lone pairs, so it can form up to four hydrogen bonds.
        \ce{HF} has one hydrogen and three lone pairs, so it is limited to
        about two. The number of hydrogen bonds per molecule matters more
        here than the strength of any single one.}
  \item State what would happen to the boiling point of \ce{CH4} if it
        were replaced by \ce{CCl4}, and justify.
        \answerlines[2]{It rises sharply. \ce{CCl4} has far more
        electrons, so its electron cloud is more polarizable and its
        London dispersion forces are much stronger, requiring more energy
        to separate the molecules.}
\end{parts}

\begin{rubric}
  \rpoint{2}{(a) 2 pts all four correct; 1 pt for three. Listing
    ``dispersion'' for \ce{HF} or \ce{H2O} as the \emph{strongest} force
    is wrong even though dispersion is present}
  \rpoint{3}{(b) 1 pt acknowledging \ce{HCl} has stronger dispersion;
    1 pt hydrogen bonding in \ce{HF}; 1 pt tied to the small, highly
    electronegative F leaving the H nearly bare}
  \rpoint{3}{(c) 1 pt water has two hydrogens and two lone pairs; 1 pt so
    it forms more hydrogen bonds per molecule; 1 pt number outweighs
    individual strength}
  \rpoint{2}{(d) 1 pt boiling point rises; 1 pt more electrons means a
    more polarizable cloud and stronger dispersion}
\end{rubric}

% =====================================================================
\frq{long}{10}{Zumdahl \S10.8 \textbullet\ CED 6.5, 3.3 \textbullet\ energy of phase changes}

A \SI{50.0}{\gram} sample of ice at \SI{0}{\celsius} is heated until it
has all become steam at \SI{100}{\celsius}.

\begin{parts}
  \item Calculate the energy needed to melt the ice.
        \answerlines[2]{$q = (50.0)(334) = \num{16700}$~J
        $= \mathbf{16.7}$~kJ}
  \item Calculate the energy needed to raise the resulting water from
        \SI{0}{\celsius} to \SI{100}{\celsius}.
        \answerlines[2]{$q = (50.0)(4.18)(100) = \num{20900}$~J
        $= \mathbf{20.9}$~kJ}
  \item Calculate the energy needed to vaporize the water at
        \SI{100}{\celsius}, and hence the total energy for the whole
        process. \workspace[2.8cm]
        \keyonly{\begin{apcanswer}$q = (50.0)(2260) = \num{113000}$~J
        $= \SI{113}{\kilo\joule}$.

        Total $= 16.7 + 20.9 + 113 =
        \mathbf{151}$~kJ\end{apcanswer}}
  \item Vaporization requires far more energy than fusion. Explain why, in
        terms of what happens to the intermolecular forces.
        \workspace[3cm]
        \keyonly{\begin{apcanswer}Melting only \emph{loosens} the
        structure. The molecules leave their fixed lattice positions and
        become free to move past one another, but they remain in contact
        and most of the hydrogen bonding survives.

        Vaporization must separate the molecules completely, so
        essentially \emph{all} the intermolecular attractions have to be
        overcome, and the molecules must be moved far apart against those
        attractions. Breaking every bond costs much more than merely
        loosening them.\end{apcanswer}}
  \item State what happens to the temperature during melting, and explain
        why. \answerlines[2]{It stays constant at \SI{0}{\celsius}. The
        energy supplied goes into overcoming intermolecular attractions
        --- increasing potential energy --- not into increasing the
        average kinetic energy, and temperature measures the latter.}
\end{parts}

\begin{rubric}
  \rpoint{2}{(a) 1 pt using $\Delta H_{\text{fus}}$ with mass only, no
    $\Delta T$; 1 pt \SI{16.7}{\kilo\joule}}
  \rpoint{2}{(b) 1 pt $mc\Delta T$ with $\Delta T = 100$; 1 pt
    \SI{20.9}{\kilo\joule}}
  \rpoint{2}{(c) 1 pt \SI{113}{\kilo\joule}; 1 pt total
    \SI{151}{\kilo\joule}}
  \rpoint{2}{(d) 1 pt melting only loosens the lattice with contact
    retained; 1 pt vaporizing separates molecules completely, overcoming
    essentially all attractions}
  \rpoint{2}{(e) 1 pt temperature constant; 1 pt energy raises potential
    energy, not kinetic energy. ``The heat is used up'' earns 0}
\end{rubric}

% =====================================================================
\frq{short}{4}{Zumdahl \S10.3--\S10.7 \textbullet\ CED 3.2 \textbullet\ types of solid}

Classify each solid and give one property that follows from your
classification.

\begin{parts}
  \item \ce{SiO2} (quartz) \answerlines[2]{\textbf{Covalent network}. Very
        high melting point, because melting requires breaking a continuous
        lattice of strong covalent bonds. Hard and non-conducting.}
  \item \ce{NaCl} \answerlines[2]{\textbf{Ionic}. High melting point and
        brittle; conducts only when molten or dissolved, when the ions
        become mobile.}
  \item \ce{Cu} \answerlines[2]{\textbf{Metallic}. Conducts electricity in
        the solid state, and is malleable, because the delocalized
        electrons move freely and tolerate displacement of the cations.}
\end{parts}

\begin{rubric}
  \rpoint{3}{1 pt each correct classification}
  \rpoint{1}{a property correctly \emph{derived} from the bonding in at
    least two cases. Listing a property without connecting it to the
    structure earns 0 for this point}
\end{rubric}

% =====================================================================
\frq{short}{4}{Zumdahl \S10.4 \textbullet\ CED 2.4 \textbullet\ metals and alloys}

\begin{parts}
  \item Explain why metals conduct electricity and are malleable, using
        the electron sea model. \workspace[2.8cm]
        \keyonly{\begin{apcanswer}A metal is a lattice of cations
        immersed in electrons that are \textbf{delocalized} over the whole
        structure rather than tied to individual atoms.

        \textbf{Conduction:} those electrons are free to drift through the
        solid when a potential difference is applied.

        \textbf{Malleability:} when a layer of cations is displaced, the
        electron sea simply flows with it and continues to bind the
        cations together. Nothing brings like charges face to face, so the
        metal deforms instead of shattering the way an ionic crystal
        does.\end{apcanswer}}
  \item Bronze is harder than pure copper. Give the structural reason.
        \answerlines[2]{Bronze is an alloy containing tin atoms of a
        different size. They disrupt the regular layers of copper atoms,
        making it harder for one layer to slide over another, so the metal
        resists deformation.}
\end{parts}

\begin{rubric}
  \rpoint{2}{(a) 1 pt delocalized electrons free to move for conduction;
    1 pt layers slide while the electron sea maintains bonding}
  \rpoint{2}{(b) 1 pt differently sized atoms disrupt the regular lattice;
    1 pt layers slide less easily}
\end{rubric}

% =====================================================================
\frq{short}{4}{Zumdahl \S10.8 \textbullet\ CED 3.3 \textbullet\ vapour pressure}

\begin{parts}
  \item Define the vapour pressure of a liquid.
        \answerlines[2]{The pressure exerted by the vapour when it is in
        dynamic equilibrium with its liquid at a given temperature ---
        molecules leaving the surface and returning at equal rates.}
  \item Diethyl ether has a much higher vapour pressure than water at
        \SI{25}{\celsius}. State what this implies about its
        intermolecular forces and its boiling point.
        \answerlines[3]{Its intermolecular forces are \textbf{weaker} ---
        it cannot hydrogen bond as water does --- so a larger fraction of
        its molecules have enough energy to escape the surface at any
        temperature. Its boiling point is correspondingly \textbf{lower},
        since less thermal energy is needed for the vapour pressure to
        reach atmospheric pressure.}
\end{parts}

\begin{rubric}
  \rpoint{2}{(a) 1 pt pressure of vapour above the liquid; 1 pt at dynamic
    equilibrium. Omitting equilibrium earns 1}
  \rpoint{2}{(b) 1 pt weaker intermolecular forces; 1 pt lower boiling
    point, linked to reaching atmospheric pressure sooner}
\end{rubric}

\examtotal

\end{document}
